\documentclass[10pt,a4paper]{article}

%double si on veut une version prof et une version élève. simple si on veut une seule version.
\newcommand{\buildmode}{simple}

\InputIfFileExists{version.tex}{}{}
% Version (définie par le script)
\providecommand{\version}{eleve}

\usepackage{feuilleinterro}

%%%à modifier à chaque fois

\renewcommand{\niveau}{1G}
\renewcommand{\chapitre}{Probabilités conditionnelles}
\renewcommand{\numchap}{2}
\renewcommand{\theme}{Probabilités}

%%%arbre pondéré à deux niveaux
%%% \arbreproba{A}{B}{P(A)}{P(Abar)}{P_A(B)}{P_A(Bbar)}{P_Abar(B)}{P_Abar(Bbar)}
\newcommand{\arbreproba}[8]{%
\begin{tikzpicture}[
  grow'=right,
  level distance=3.6cm,
  level 1/.style={sibling distance=2.5cm},
  level 2/.style={sibling distance=1.25cm},
  edge from parent/.style={draw,-latex},
  every node/.style={font=\small}
]
\node {}
  child { node {\(#1\)}
    child { node {\(#2\)} edge from parent node[above] {#5} }
    child { node {\(\overline{#2}\)} edge from parent node[below] {#6} }
    edge from parent node[above] {#3}
  }
  child { node {\(\overline{#1}\)}
    child { node {\(#2\)} edge from parent node[above] {#7} }
    child { node {\(\overline{#2}\)} edge from parent node[below] {#8} }
    edge from parent node[below] {#4}
  };
\end{tikzpicture}%
}
\newcommand{\vide}{\makebox[1.3cm]{\dotfill}}

%%%cases du tableau plus hautes pour que les élèves puissent écrire dedans
\renewcommand{\arraystretch}{1.8}
\newcolumntype{C}{>{\centering\arraybackslash}p{2.4cm}}

\begin{document}

% =========================================================
% PAGE 1 : VERSION A
% =========================================================

\interroversion{A}

\textbf{Exercice 1 -- Calculs avec des fractions \hfill /2}

Pour \(x\neq2\), écrire l'expression suivante sous la forme d'une seule fraction, en détaillant les étapes :
\[
A=\dfrac{2}{3}+\dfrac{x+1}{x-2}.
\]

\vspace{2.3cm}

\textbf{Exercice 2 -- Demi-pension et sport \hfill /8}

Un lycée compte \(200\) élèves de première. On sait que :
\begin{itemize}
\item \(60\,\%\) des élèves sont demi-pensionnaires, les autres sont externes ;
\item \(31\,\%\) des élèves pratiquent un sport en club.
\end{itemize}

\begin{enumerate}
\item Compléter le tableau ci-dessous, en justifiant les calculs utilisant les pourcentages. \hfill /2,5

\begin{center}
\begin{tabular}{|l|C|C|C|}
    \hline
     & Sport en club & Pas de sport & Total \\
    \hline
    Demi-pensionnaires & & 90 & \\
    \hline
    Externes & 32 & & \\
    \hline
    Total & & & 200 \\
    \hline
\end{tabular}
\end{center}

\vspace{0.8cm}

On choisit un élève de première au hasard. On note \(D\) : \og l'élève est demi-pensionnaire \fg{} et \(S\) : \og l'élève pratique un sport en club \fg{}.

\item Calculer \(P(D)\). \hfill /1

\vspace{1.1cm}

\item Calculer \(P(D\cap S)\). \hfill /1

\vspace{1.1cm}

\item Calculer \(P_D(S)\). \hfill /1

\vspace{1.1cm}

\item À l'aide du tableau, compléter l'arbre pondéré ci-dessous. \hfill /1,5

\begin{center}
\arbreproba{D}{S}{\vide}{\vide}{\vide}{\vide}{\vide}{\vide}
\end{center}

\item À l'aide de l'arbre, retrouver \(P(D\cap S)\) en rappelant la formule utilisée. \hfill /1

\vspace{2cm}
\end{enumerate}

\newpage

% =========================================================
% PAGE 2 : VERSION B
% =========================================================

\interroversion{B}

\textbf{Exercice 1 -- Calculs avec des fractions \hfill /2}

Pour \(x\neq-3\), écrire l'expression suivante sous la forme d'une seule fraction, en détaillant les étapes :
\[
A=\dfrac{3}{4}+\dfrac{x-1}{x+3}.
\]

\vspace{2.3cm}

\textbf{Exercice 2 -- Au cinéma \hfill /8}

Un soir, un cinéma a accueilli \(400\) spectateurs. On sait que :
\begin{itemize}
\item \(40\,\%\) des spectateurs ont payé un tarif réduit, les autres ont payé plein tarif ;
\item \(27\,\%\) des spectateurs ont vu un film en 3D, les autres un film en 2D.
\end{itemize}

\begin{enumerate}
\item Compléter le tableau ci-dessous, en justifiant les calculs utilisant les pourcentages. \hfill /2,5

\begin{center}
\begin{tabular}{|l|C|C|C|}
    \hline
     & Film en 3D & Film en 2D & Total \\
    \hline
    Tarif réduit & & 112 & \\
    \hline
    Plein tarif & 60 & & \\
    \hline
    Total & & & 400 \\
    \hline
\end{tabular}
\end{center}

\vspace{0.8cm}

On choisit un spectateur au hasard. On note \(R\) : \og le spectateur a payé un tarif réduit \fg{} et \(T\) : \og le spectateur a vu un film en 3D \fg{}.

\item Calculer \(P(R)\). \hfill /1

\vspace{1.1cm}

\item Calculer \(P(R\cap T)\). \hfill /1

\vspace{1.1cm}

\item Calculer \(P_R(T)\). \hfill /1

\vspace{1.1cm}

\item À l'aide du tableau, compléter l'arbre pondéré ci-dessous. \hfill /1,5

\begin{center}
\arbreproba{R}{T}{\vide}{\vide}{\vide}{\vide}{\vide}{\vide}
\end{center}

\item À l'aide de l'arbre, retrouver \(P(R\cap T)\) en rappelant la formule utilisée. \hfill /1

\vspace{2cm}
\end{enumerate}

\newpage

% =========================================================
% PAGE 3 : VERSION C
% =========================================================

\interroversion{C}

\textbf{Exercice 1 -- Calculs avec des fractions \hfill /2}

Pour \(x\neq1\), écrire l'expression suivante sous la forme d'une seule fraction, en détaillant les étapes :
\[
A=\dfrac{1}{5}+\dfrac{2x+1}{x-1}.
\]

\vspace{2.3cm}

\textbf{Exercice 2 -- Clients d'un magasin \hfill /8}

Une journée, un magasin a reçu \(500\) clients. On sait que :
\begin{itemize}
\item \(60\,\%\) des clients possèdent la carte de fidélité du magasin ;
\item \(30\,\%\) des clients ont fait un achat de plus de \(50\) \euro.
\end{itemize}

\begin{enumerate}
\item Compléter le tableau ci-dessous, en justifiant les calculs utilisant les pourcentages. \hfill /2,5

\begin{center}
\begin{tabular}{|l|C|C|C|}
    \hline
     & Plus de 50 \euro & 50 \euro{} ou moins & Total \\
    \hline
    Avec carte & & 180 & \\
    \hline
    Sans carte & 30 & & \\
    \hline
    Total & & & 500 \\
    \hline
\end{tabular}
\end{center}

\vspace{0.8cm}

On choisit un client au hasard. On note \(F\) : \og le client possède la carte de fidélité \fg{} et \(M\) : \og le client a fait un achat de plus de \(50\) \euro{} \fg{}.

\item Calculer \(P(F)\). \hfill /1

\vspace{1.1cm}

\item Calculer \(P(F\cap M)\). \hfill /1

\vspace{1.1cm}

\item Calculer \(P_F(M)\). \hfill /1

\vspace{1.1cm}

\item À l'aide du tableau, compléter l'arbre pondéré ci-dessous. \hfill /1,5

\begin{center}
\arbreproba{F}{M}{\vide}{\vide}{\vide}{\vide}{\vide}{\vide}
\end{center}

\item À l'aide de l'arbre, retrouver \(P(F\cap M)\) en rappelant la formule utilisée. \hfill /1

\vspace{2cm}
\end{enumerate}

\newpage

% =========================================================
% PAGE 4 : VERSION D
% =========================================================

\interroversion{D}

\textbf{Exercice 1 -- Calculs avec des fractions \hfill /2}

Pour \(x\neq-1\), écrire l'expression suivante sous la forme d'une seule fraction, en détaillant les étapes :
\[
A=\dfrac{3}{2}+\dfrac{x+2}{x+1}.
\]

\vspace{2.3cm}

\textbf{Exercice 2 -- Une course à pied \hfill /8}

\(250\) coureurs ont participé à une course à pied. On sait que :
\begin{itemize}
\item \(40\,\%\) des coureurs sont licenciés dans un club d'athlétisme ;
\item \(36\,\%\) des coureurs ont terminé la course en moins d'une heure.
\end{itemize}

\begin{enumerate}
\item Compléter le tableau ci-dessous, en justifiant les calculs utilisant les pourcentages. \hfill /2,5

\begin{center}
\begin{tabular}{|l|C|C|C|}
    \hline
     & Moins d'1 h & 1 h ou plus & Total \\
    \hline
    Licenciés & & 40 & \\
    \hline
    Non licenciés & 30 & & \\
    \hline
    Total & & & 250 \\
    \hline
\end{tabular}
\end{center}

\vspace{0.8cm}

On choisit un coureur au hasard. On note \(L\) : \og le coureur est licencié \fg{} et \(H\) : \og le coureur a terminé en moins d'une heure \fg{}.

\item Calculer \(P(L)\). \hfill /1

\vspace{1.1cm}

\item Calculer \(P(L\cap H)\). \hfill /1

\vspace{1.1cm}

\item Calculer \(P_L(H)\). \hfill /1

\vspace{1.1cm}

\item À l'aide du tableau, compléter l'arbre pondéré ci-dessous. \hfill /1,5

\begin{center}
\arbreproba{L}{H}{\vide}{\vide}{\vide}{\vide}{\vide}{\vide}
\end{center}

\item À l'aide de l'arbre, retrouver \(P(L\cap H)\) en rappelant la formule utilisée. \hfill /1

\vspace{2cm}
\end{enumerate}

\end{document}
